\documentclass[11pt,a4paper]{article}

\usepackage{amsmath,amssymb,amsthm}
\usepackage{geometry}
\usepackage{hyperref}
\usepackage{booktabs}
\usepackage{array}
\usepackage{enumitem}
\usepackage{titlesec}
\usepackage{xcolor}
\usepackage{fancyhdr}
\usepackage[expansion=false]{microtype}
\usepackage{parskip}
\usepackage{listings}
\usepackage{mdframed}

\geometry{margin=2.5cm}

\definecolor{deepblue}{RGB}{31,56,100}
\definecolor{midblue}{RGB}{46,84,150}
\definecolor{codebg}{RGB}{245,247,250}
\definecolor{coderule}{RGB}{200,210,225}
\definecolor{warnred}{RGB}{160,40,40}

\titleformat{\section}{\large\bfseries\color{deepblue}}{{\thesection}}{1em}{}
\titleformat{\subsection}{\normalsize\bfseries\color{midblue}}{{\thesubsection}}{1em}{}
\titleformat{\subsubsection}{\normalsize\bfseries}{{\thesubsubsection}}{1em}{}

\lstset{
  backgroundcolor=\color{codebg},
  basicstyle=\ttfamily\small,
  breaklines=true,
  frame=single,
  rulecolor=\color{coderule},
  framesep=6pt,
  xleftmargin=8pt,
  xrightmargin=8pt,
  aboveskip=8pt,
  belowskip=8pt,
}

\newmdenv[
  backgroundcolor=codebg,
  linecolor=coderule,
  linewidth=0.5pt,
  innerleftmargin=10pt,
  innerrightmargin=10pt,
  innertopmargin=8pt,
  innerbottommargin=8pt,
  skipabove=8pt,
  skipbelow=8pt,
]{shaded}

\pagestyle{fancy}
\fancyhf{}
\fancyhead[L]{\small\color{gray}Linearized GR Sandbox}
\fancyhead[R]{\small\color{gray}v38 Addendum\,---\,Field Initialization \& Absorbing BCs}
\fancyfoot[C]{\small\thepage}
\renewcommand{\headrulewidth}{0.4pt}

\newcommand{\vv}[1]{\mathbf{#1}}
\newcommand{\Phig}{\Phi_g}
\newcommand{\Ag}{\vv{A}_g}
\newcommand{\order}[1]{\mathcal{O}(#1)}
\newcommand{\dd}{\mathrm{d}}
\newcommand{\half}{\tfrac{1}{2}}

% Numbered algorithm steps (mirrors gr_sandbox_v38.tex).
\newlist{algsteps}{enumerate}{1}
\setlist[algsteps]{label=\textbf{Step \arabic*.}, leftmargin=*,
                   labelwidth=5em, itemsep=6pt}

%--------------------------------------------------------------------
\begin{document}
%--------------------------------------------------------------------

\begin{center}
  {\Large\bfseries\color{deepblue} v38 Addendum: Field Initialization and \\[2pt]
   Absorbing Boundary Conditions Revisited\par}
  \vspace{0.4cm}
  {\small Linearized GR + EM PIC Sandbox, web-frontend phase\par}
  \vspace{0.2cm}
  {\small \today\par}
\end{center}

\vspace{0.5cm}

\begin{shaded}
\textbf{Scope.}  This addendum revises three sections of
gr\_sandbox\_v38.tex that, while internally consistent, were found to
be physically wrong or numerically broken in the course of building
the interactive web demo against the v40/v41 simulator core:
\begin{enumerate}[leftmargin=*,nosep]
  \item \S 9.6 ``Absorbing boundary layer'': the multiplicative
    ($\sigma\,\dd t$) damping ring does not absorb the lowest standing
    mode of the bounded box.  Replace with a derivative-friction
    (Klein-Gordon-like) absorber.
  \item \S 15.9 ``Field initialization'': the wave-equation
    convergence iteration starting from $\Phi=0$ does not converge to
    a static solution for the same reason -- the lowest standing
    mode, seeded by the initial transient, is undamped.  Replace
    with a direct-sum 2D Li{\'e}nard-Wiechert initialization that
    sets all six potentials analytically.
  \item Add a new subsection on the DC gauge mode: under Neumann
    outer BC the constant-everywhere mode of a scalar potential is
    accelerated by any nonzero source mean and runs away; under
    Dirichlet the wall pins it.  This was implicit in v38 but the
    failure mode wasn't called out.
\end{enumerate}
The discrete stability analysis of the centered-implicit friction
form (which the post-step ``bleed'' form is \emph{not}) is included
because the result is non-trivial and easy to get wrong.
\end{shaded}

\begin{shaded}
\textbf{Revision 2 (2026-06-05).}  Corrections from running the demo
end-to-end, marked inline as ``Revision 2'' boxes and collected in
\S\ref{sec:status}:
\begin{enumerate}[leftmargin=*,nosep]
  \item Field init uses a per-particle \emph{boundary-mean subtraction}
    (a gauge shift) plus a \emph{frozen-friction settle}, not the
    multiplicative taper of the first draft (which distorts gradients).
  \item The friction absorber uses a \emph{zero interior floor}
    (wall-ramp only) for dynamics: a nonzero floor damps the
    time-varying orbital near-field of moving sources and causes a
    numerical inspiral.
  \item The pic\_binary COM drift was a \emph{TSC deposition} artifact,
    fixed by the \emph{BUMP $R=8$} kernel -- not, as first conjectured,
    the init transient.  The near-circular orbit also needs the
    discrete-force velocity calibration $v_\text{factor}\approx 1.05$.
\end{enumerate}
\end{shaded}

\tableofcontents

\vspace{0.5cm}

%--------------------------------------------------------------------
\section{Diagnosis of the v38 absorber}
%--------------------------------------------------------------------

The v38 absorbing ring multiplies the leapfrog update by
$(1 - \sigma_k\,\dd t)$ in a polynomial-tapered ring of depth $N$
cells:
\begin{equation}
  \Phi_k^{n+1} \;=\; \bigl[2\Phi_k^n - \Phi_k^{n-1}
                       + c^2\dd t^2\,\bigl(\nabla^2\Phi^n + s_c\,\rho_k\bigr)\bigr]
                  \,\bigl(1 - \sigma_k\,\dd t\bigr).
\end{equation}
This factor does two distinct things: (a) it dissipates outgoing wave
energy, which is what we want; and (b) it pulls $\Phi$ toward zero
everywhere inside the ring, which effectively imposes a Dirichlet
boundary at the wall.  Effect (a) only acts on field amplitude that
\emph{reaches} the ring; effect (b) is the reason the static
fixed point of the box is the Dirichlet Poisson solution rather
than the free-space one.

The failure mode is mode-dependent.  Consider the standing modes
$\phi_{mn}(x,y) = \sin(m\pi x/L)\sin(n\pi y/L)$ of a square box with
zero-Dirichlet walls.  The lowest mode $(m,n) = (1,1)$ has its
\textbf{antinode at the center} and \textbf{nodes at the wall}.  Its
overlap with the multiplicative damping array, concentrated in the
outer ring, is suppressed by the eigenfunction's small wall
amplitude.  Specifically, the per-step decay factor for mode
$\phi_{mn}$ is approximately
\begin{equation}
  1 - \dd t\,\overline{\sigma}_{mn},
  \qquad
  \overline{\sigma}_{mn}
     \;=\;
     \frac{\int_\Omega \sigma(x,y)\,\phi_{mn}^2\,\dd^2 x}
          {\int_\Omega \phi_{mn}^2\,\dd^2 x}.
\end{equation}
For the lowest mode this average is $\order{N/L}$ times the wall value
of $\sigma$; for high-$k$ modes (which oscillate with appreciable
amplitude at the wall) it is $\order{1}$ times the wall value.  At
$N/L \sim 1/8$ (a typical absorber ring depth) the lowest mode is
damped roughly an order of magnitude slower than high-$k$ modes.

\textbf{Empirical signature.}  With pic\_static (a single stationary
particle), particles frozen, and the field initialized by the v38
\S 15.9 convergence iteration, the value of $\Phig$ at the box center
oscillates with period $\sim 2L/c$ and amplitude $\sim |\Phig^\text{static}|$
indefinitely.  The probe trajectory shows envelope swing of
$\sim 50\%$ of $|\Phig^\text{static}|$ after 4000 steps with no decay
trend.  The absorbing ring is not absorbing the lowest mode.

\vspace{0.5em}
The v38 \S 15.9 prescription -- ``iterate the leapfrog with frozen
particles until the field converges'' -- relies on the absorber to
remove the transient launched by switching the source on at $t=0$.
That transient has Fourier content at every mode; the lowest
component is excited as strongly as any other, and once excited it
persists.  The convergence iteration therefore does not, in fact,
converge to the static solution; it converges to the static solution
\emph{plus} a residual standing-mode oscillation of comparable
amplitude.

%--------------------------------------------------------------------
\section{Direct-sum Liénard-Wiechert initialization}
%--------------------------------------------------------------------
\label{sec:lw_init}

The correct initialization is to set each perturbation potential at
$t=0$ to the static Poisson solution of the deposited sources, then
start the wave equation from there with zero time-derivative.  For an
isolated $N$-particle configuration in 2D, the static solution is
known analytically -- it is the 2D Liénard-Wiechert potential, which
under the Lorenz gauge is built from the same $\ln r$ kernel for all
six fields:
\begin{align}
  \Phig(\vv{x})  &\;=\; -\frac{s_c^g}{2\pi}\,\sum_i m_i\,\ln|\vv{x}-\vv{x}_i|,
    \label{eq:phig_lw} \\
  \Ag(\vv{x})  &\;=\; -\frac{s_c^g}{2\pi c^2}\,\sum_i m_i\,\vv{v}_i\,\ln|\vv{x}-\vv{x}_i|,
    \label{eq:Ag_lw} \\
  \phi(\vv{x})  &\;=\; -\frac{s_c^{em}}{2\pi}\,\sum_i q_i\,\ln|\vv{x}-\vv{x}_i|,
    \label{eq:phi_lw} \\
  \vv{A}(\vv{x})  &\;=\; -\frac{s_c^{em}}{2\pi c^2}\,\sum_i q_i\,\vv{v}_i\,\ln|\vv{x}-\vv{x}_i|,
    \label{eq:A_lw}
\end{align}
where $s_c^g = -4\pi G_{\text{eff}}$ and $s_c^{em} = +4\pi k_e$ are
the source coefficients from \S 9.2, and $\vv{v}_i = \vv{p}_i/(\gamma_i
m_i)$.  These follow from the 2D fundamental solution of $\nabla^2$,
$G(r) = (2\pi)^{-1}\ln r$, applied to point-source charge/mass and
current distributions $\rho(\vv{y}) = \sum_i Q_i\,\delta(\vv{y}-\vv{x}_i)$.
The velocity moves from the scalar to the vector exactly as in the
3D Liénard-Wiechert form; in the static-source limit the scalar
parts dominate and the vector parts are an $\order{v/c}$ correction.

\subsubsection*{Sub-cell clip}

The logarithm diverges at $\vv{x} = \vv{x}_i$.  Clip the radius at
$r_{\min} = \half\,\dd x$ (smaller than any cell-particle separation
that can actually occur on the half-cell-centered deposition lattice)
when evaluating the analytic sum; the resulting value at the source
cell is finite and within the floating-point dynamic range of all
nearby cells.  The clipped value is not the discrete 5-point Poisson
solution -- see \S\ref{sec:lw_discrete} -- but the residual mismatch is
handled by the friction absorber introduced below.

\subsubsection*{Boundary-mean subtraction (supersedes the taper)}

\begin{shaded}
\textbf{Revision 2 (2026-06-05).}  An earlier draft of this addendum
multiplied the L-W field by a smooth taper $T(d)$ to bring it to zero
through the absorbing zone.  That taper is \emph{not} gauge-clean --
multiplying the field distorts its gradient (hence the force) inside
the ring -- and it is no longer used.  It is replaced by the
boundary-mean subtraction below.
\end{shaded}

Under zero-Dirichlet the bare L-W field is large and \emph{positive} at
the wall: the 2D log \emph{grows} outward, so $\Phi(\text{wall}) \sim
-(s_c/2\pi)\,m\,\ln R_\text{wall}$, which for the binary exceeds the
near-field well depth by a factor $\sim 7$--$15$.  Forcing that to zero
is the dominant startup transient.

The fix is a pure \emph{gauge} shift: for each particle $i$, subtract
the mean of $\ln r$ over the boundary ring from that particle's
contribution before accumulating it,
\begin{equation}
  \ln|\vv{x}-\vv{x}_i| \;\longrightarrow\;
  \ln|\vv{x}-\vv{x}_i| \;-\; \big\langle \ln|\vv{x}_b - \vv{x}_i|\big\rangle_{b\,\in\,\partial\Omega},
\end{equation}
applied uniformly to all six fields (the scalar shift maps to $\Ag,
\vv{A}$ through the same $\vv{v}/c^2$ factor).  Because it is a constant
per particle, it changes no gradient and hence no force anywhere.  The
summed field then has zero mean on the boundary ring -- consistent with
the Dirichlet fixed point -- removing $\sim 93\%$ of the boundary
mismatch (the residual is the $\sim 7\%$ spatial variation of $\ln r$
around the ring, $\ln\sqrt2 / \ln R_\text{wall}$, which a constant
cannot remove).  Under Neumann BC the subtraction is harmless but
unnecessary.

\subsubsection*{Frozen-particle friction settle}

The boundary-mean L-W is a close \emph{initial guess} but still differs
from the discrete fixed point by the near-source discrete-vs-continuum
mismatch (\S\ref{sec:lw_discrete}, $\sim 4\%$ peak).  Relax it to the
fixed point by iterating the field with the particles \emph{frozen} and
the friction absorber on (\texttt{gr\_sim\_init\_potentials\_settled}):
positions and velocities are pinned, so the sources are constant and
the wave equation -- now genuinely dissipative thanks to the friction
of \S\ref{sec:friction} -- converges to the static solution.  This is
the v38 \S 15.9 convergence iteration, which \emph{only works now} that
friction damps the lowest mode (\S\ref{sec:friction}); the close L-W
guess makes it converge in a few hundred steps.

Measured (256$\times$256 \texttt{pic\_binary}): the symmetric binary's
quarter-orbit total-momentum impulse falls from $2.3\times10^{-3}$ of
$m v_\text{orb}$ with no settle to $\sim 3\times10^{-5}$ (round-off)
with $\gtrsim 100$ settle steps.  Wall-ramp-only friction (no interior
floor) settles equally well, so the settle does \emph{not} require an
interior friction floor -- see the caution in \S\ref{sec:friction}.

\subsubsection*{Leapfrog initialization at $t=0$}

After the L-W direct sum is written into the curr buffer of each
field, also write it into the prev and next buffers:
\begin{equation}
  \Phi^{-1}_k \;=\; \Phi^{0}_k \;=\; \Phi^{1}_k \;=\; \Phi^\text{LW}_k.
\end{equation}
This makes the discrete time derivative encoded by the leapfrog
($\Phi^0 - \Phi^{-1}$) exactly zero, so no spurious transient is
launched by the time-integration startup.  Particle momenta are set
directly from the scenario initial conditions; the v38 \S 15.9
half-step back-extrapolation is not needed because the field is
already at the static solution.

\subsubsection*{Compute cost}

For $N_p$ particles on a $W \times H$ grid the L-W sum is
$\order{N_p \cdot W \cdot H}$ log evaluations.  At $W=H=384$, $N_p=2$
this is $\sim 3 \times 10^5$ operations -- sub-millisecond.  For
$N_p \gtrsim 10^4$ a tree-code or FFT convolution would be needed,
but the regime the sandbox addresses (a few orbiting test masses)
does not require it.

%--------------------------------------------------------------------
\section{Continuum--discrete Poisson mismatch}
%--------------------------------------------------------------------
\label{sec:lw_discrete}

The L-W direct sum is the analytic Poisson solution of a sum of Dirac
delta sources.  The simulation's actual source is not a delta; it is
the TSC-spread (or BUMP, or CIC) deposition kernel evaluated at the
particle's sub-cell position.  Inserting the analytic L-W into the
discrete 5-point leapfrog and evaluating the residual
$\nabla^2_5 \Phi^\text{LW} + s_c\,\rho^\text{TSC}$ shows that the two
are not equal at the source.  At a single particle of mass $m$
centered exactly on a cell, with TSC weights $(9/16, 3/32, 1/64)$ on
the $3 \times 3$ stamp:
\begin{equation}
  \nabla^2_5 \Phi^\text{LW}\bigr|_\text{center}
  \;=\;
  -\frac{4\,\Phi^\text{LW}_\text{center}}{\dd x^2}
  \;\approx\; +5.55 \times 10^{-2}
  \qquad\text{(at $m = 0.01$, $G_\text{eff}=1$)},
\end{equation}
while
\begin{equation}
  s_c\,\rho^\text{TSC}_\text{center}
  \;=\; -4\pi G_\text{eff}\,\frac{m \cdot 9/16}{\dd x^2}
  \;\approx\; -7.07 \times 10^{-2}.
\end{equation}
The residual $-1.5\times 10^{-2}$ is roughly $20\%$ of $|s_c\rho|$ at
the source cell.  It produces a per-step kick on $\Phig$ at the
source of $c^2\dd t^2 \times (-1.5\times 10^{-2}) \approx -8\times
10^{-3}$ -- $\sim 5\%$ of $|\Phig^\text{LW}_\text{center}|$ every
step.  This is the source-term forcing that the friction absorber
(next section) must dissipate.

\paragraph{Why not solve the discrete 5-point Poisson exactly?}  An
SOR sweep on $\nabla^2_5 \Phi = -s_c\,\rho^\text{TSC}$ converges in
$\order{\max(W,H)}$ iterations to the bit-exact discrete Poisson
solution and removes this residual entirely.  This was implemented
and tested.  However, the SOR solution requires a boundary
condition: zero-Dirichlet at the outer ring is the natural choice,
but the resulting static field is the Dirichlet-Poisson solution of
the box, not the free-space 2D log.  The two differ visibly near
the wall, where the Dirichlet enforcement forces $\Phi \to 0$ rather
than continuing the $\ln r$ tail.  For the pedagogical sandbox the
free-space appearance of L-W is preferable, so we accept the small
discrete-vs-continuum mismatch and rely on the friction absorber to
dissipate the resulting per-step forcing.

%--------------------------------------------------------------------
\section{Derivative-friction absorber}
%--------------------------------------------------------------------
\label{sec:friction}

Replace the multiplicative damping ring of \S 9.6 with a derivative-
friction term added to the wave equation:
\begin{equation}
  \frac{\partial^2 \Phi}{\partial t^2} \;+\; 2\gamma(\vv{x})\,\frac{\partial \Phi}{\partial t}
  \;=\; c^2\nabla^2 \Phi \;+\; s_c\,\rho.
  \label{eq:friction_pde}
\end{equation}
The crucial difference from the multiplicative form is that the
friction term vanishes identically at $\partial_t \Phi = 0$, so the
static fixed point is exactly the Poisson solution regardless of the
spatial profile of $\gamma$.  In particular: a uniform friction floor
$\gamma_0$ everywhere in the box damps all modes -- including the
lowest standing mode, which has nonzero $\partial_t \Phi$ even though
its spatial amplitude at the wall is small -- without shifting the
\emph{static} equilibrium.

\emph{Caveat (Revision 2):} ``damps all modes with nonzero
$\partial_t\Phi$'' includes the legitimate, time-varying near-field of
\emph{moving} sources.  A uniform interior floor therefore bleeds real
orbital energy (numerical inspiral), so for dynamics the floor is set
to zero and the absorber is wall-ramp-only; see \S\ref{sec:friction}
``Per-cell profile''.

\subsection{Mode-by-mode decay rate}

Consider a single eigenmode $\Phi_k(\vv{x},t) = a_k(t)\phi_k(\vv{x})$
of the spatial operator $-c^2\nabla^2$ with eigenvalue $\omega_k^2$,
and let $\overline{2\gamma}_k$ denote the eigenmode-averaged
friction.  In the source-free regime:
\begin{equation}
  \ddot a_k + \overline{2\gamma}_k\,\dot a_k + \omega_k^2\,a_k \;=\; 0
  \quad\Longrightarrow\quad
  a_k(t) \;\propto\; e^{-\overline{\gamma}_k\,t}\,\cos(\sqrt{\omega_k^2 - \overline{\gamma}_k^2}\,t + \delta).
\end{equation}
For $\overline{\gamma}_k < \omega_k$ (the underdamped regime) every
mode decays at rate $\overline{\gamma}_k$.  The lowest mode's
eigenmode-averaged friction is
$\overline{\gamma}_{(1,1)} \sim \gamma_0$ (dominated by the uniform
floor, since the wall ramp contributes only at the small-amplitude
boundary).  A uniform $\gamma_0\,\dd t = 10^{-3}$ gives mode-amplitude
decay to $1/e$ in $\sim 1000$ steps.  Higher-$k$ modes are also
damped by the wall ramp; combined decay rates approach $\gamma_\text{wall}$
for modes localized near the boundary.

\subsection{Discrete leapfrog: post-step bleed (UNSTABLE)}

The natural-looking discretization of (\ref{eq:friction_pde}) is to
run the standard leapfrog and then, after each rotation, bleed the
prev buffer toward curr:
\begin{equation}
  \Phi^{n-1}_{k,\text{bled}} \;:=\; (1-\alpha_k)\,\Phi^{n-1}_k \;+\; \alpha_k\,\Phi^n_k,
  \qquad \alpha_k = 2\gamma_k\,\dd t.
  \label{eq:post_bleed}
\end{equation}
This appears to reduce the encoded time derivative $\Phi^n - \Phi^{n-1}$
by a factor $(1-\alpha_k)$ each step without touching the static
fixed point (where $\Phi^n = \Phi^{n-1}$ and the bleed is a no-op).

\textbf{This form is unconditionally unstable at the 2D Nyquist mode
for any $\alpha > 0$ when CFL $= 1/\sqrt{2}$.}

The recurrence implied by (\ref{eq:post_bleed}) is
\begin{equation}
  \Phi^{n+2} \;=\; (2 - \alpha - s^2)\,\Phi^{n+1} \;-\; (1-\alpha)\,\Phi^n
  \;+\; c^2\dd t^2 \cdot \text{(source terms)},
\end{equation}
where $s^2 = c^2 \dd t^2 |\nabla^2_5|$.  At the 2D Nyquist
($s^2 = 4$ when CFL $= 1/\sqrt{2}$), the characteristic equation
$\lambda^2 - (2-\alpha-s^2)\lambda + (1-\alpha) = 0$ has discriminant
$\alpha^2 + 8\alpha > 0$ for any $\alpha > 0$.  The roots are
\begin{equation}
  \lambda_{\pm} \;=\; \frac{-2 - \alpha \;\pm\; \sqrt{\alpha^2 + 8\alpha}}{2}.
\end{equation}
For $\alpha = 10^{-3}$, $\lambda_- \approx -1.045$ -- the Nyquist
mode grows $4.5\%$ per step.  Round-off seeds it; saturation in
$\sim 100$ steps.

Empirically (the failure mode that prompted this addendum): on a
$384 \times 384$ grid with L-W initialization and post-step-bleed
friction at $\alpha = 0.04$ near the wall, a coherent wavefront
appears at the boundary within a few hundred steps, propagates
inward, and dominates the entire field within a few thousand.  The
``absorbing layer'' is acting as a generator.

\subsection{Centered-implicit discretization (STABLE)}

Discretize the friction term at the leapfrog midpoint, so the friction
gradient term $\partial_t \Phi$ uses $(\Phi^{n+1} - \Phi^{n-1})/(2\dd
t)$:
\begin{equation}
  \frac{\Phi^{n+1} - 2\Phi^n + \Phi^{n-1}}{\dd t^2}
  \;+\;
  2\gamma\,\frac{\Phi^{n+1} - \Phi^{n-1}}{2\,\dd t}
  \;=\;
  c^2\nabla^2_5 \Phi^n \;+\; s_c\,\rho^n.
\end{equation}
Solving for $\Phi^{n+1}$ with $\beta = \gamma\,\dd t$:
\begin{equation}
  (1 + \beta)\,\Phi^{n+1} \;=\; 2\Phi^n \;-\; (1-\beta)\,\Phi^{n-1}
  \;+\; \dd t^2\,\bigl(c^2\nabla^2_5 \Phi^n + s_c\,\rho^n\bigr).
  \label{eq:centered}
\end{equation}
At the static fixed point $\Phi^{n+1} = \Phi^n = \Phi^{n-1}$ the
LHS factor $(1+\beta)$ and the RHS coefficient $-(1-\beta) + 2 = 1+\beta$
match identically -- the static is preserved.

\paragraph{Stability at 2D Nyquist.}  The characteristic equation
\begin{equation}
  (1+\beta)\lambda^2 \;-\; (2 - s^2)\lambda \;+\; (1-\beta) \;=\; 0
\end{equation}
at $s^2 = 4$ has discriminant
$(2 - 4)^2 - 4(1-\beta)(1+\beta) = 4\beta^2$, giving real roots
\begin{equation}
  \lambda_+ \;=\; -\frac{1 - \beta}{1 + \beta},
  \qquad
  \lambda_- \;=\; -1.
\end{equation}
$|\lambda_+| < 1$ for $\beta > 0$; $|\lambda_-| = 1$.  The Nyquist
mode is marginally stable on one branch and damped on the other --
unconditionally non-growing for any $\beta \ge 0$.  At lower-$k$ modes
the discriminant is negative and $|\lambda|^2 = (1-\beta)/(1+\beta) <
1$ for $\beta > 0$: every interior mode strictly decays.

\paragraph{Post-leapfrog correction implementation.}  Equation
(\ref{eq:centered}) is equivalent to running the standard (un-friction)
leapfrog and then applying a per-cell post-correction \emph{before}
the buffer rotation:
\begin{equation}
  \Phi^{n+1}_k \;\leftarrow\; \frac{\Phi^{n+1}_{k,\text{std}}
                                    \;+\; \beta_k\,\Phi^{n-1}_k}{1 + \beta_k}.
  \label{eq:centered_post}
\end{equation}
This is the form actually used in the implementation: no leapfrog
kernel needs to be modified, only a single $\order{N_\text{cells}}$
sweep per field per step.  At $\beta_k = 0$ the formula reduces to
$\Phi^{n+1}_{k,\text{std}}$ identically, preserving backward
compatibility for tests that do not use the friction array.

\subsection{Per-cell profile}

The implementation stores $\beta_k = \gamma_k\,\dd t$ as a single
floating-point array \texttt{friction\_d[k]}.  The profile is a
polynomial ramp from an interior floor to a larger wall value:
\begin{equation}
  \beta_k \;=\;
  \begin{cases}
    \beta_\text{floor} & d_k \ge N_\text{damp} \\
    \beta_\text{floor} + (\beta_\text{wall} - \beta_\text{floor})\,u^2,\; u = (N_\text{damp} - d_k)/N_\text{damp} & 0 \le d_k < N_\text{damp}
  \end{cases}
\end{equation}
with $d_k$ the cell-count distance from the nearest wall.

\begin{shaded}
\textbf{Revision 2 (2026-06-05): use $\beta_\text{floor} = 0$ (wall-ramp
only) for dynamics.}  A nonzero interior floor was first chosen to damp
the lowest standing mode everywhere.  But friction damps \emph{every}
mode with nonzero $\partial_t\Phi$ -- including the legitimate,
time-varying orbital near-field of moving sources.  A floor of $10^{-3}$
bleeds the binary's orbital energy: a measurable numerical
\emph{inspiral} ($\sim 13\%$ orbit shrink over 4 orbits in
\texttt{pic\_binary}).  The interior floor is moreover unnecessary: with
the boundary-mean L-W initial guess, the frozen-particle settle reaches
the fixed point using the \emph{wall ramp alone} (the near-source
residual radiates to the wall), and during dynamics there is no startup
transient to trap.  The recommended values for the web demo
($384\times384$, $N_\text{damp}=48$) are therefore
\begin{equation}
  \beta_\text{floor} \;=\; 0, \qquad \beta_\text{wall} \;=\; 2\times10^{-2}.
\end{equation}
A small floor ($\sim 10^{-4}$) is acceptable if a particular scenario
excites a persistent interior mode, at the cost of a correspondingly
small inspiral; the default is $0$.
\end{shaded}

The wall ramp gives radiative wave packets several e-folds of
attenuation across the ring; modes localized near the boundary decay in
$\sim 50$ steps.  At CFL $= 1/\sqrt{2}$ this is well inside the
stability region (\S\ref{sec:friction}, any $\beta \ge 0$).

%--------------------------------------------------------------------
\section{The DC mode under each boundary condition}
%--------------------------------------------------------------------
\label{sec:dc_mode}

The friction absorber damps every mode with nonzero
$\partial_t \Phi$.  The constant-everywhere mode $\Phi_{DC}(\vv{x},t) =
a(t)$ has $\partial_t a$ that is \emph{not} zero in general -- the
wave equation drives it via the spatial mean of the source:
\begin{equation}
  \ddot a \;=\; s_c\,\bar\rho,
  \qquad
  \bar\rho \;=\; \frac{1}{|\Omega|}\int_\Omega \rho\,\dd^2 x.
\end{equation}
For a net-nonzero source $\bar\rho \ne 0$ (any scenario with nonzero
total mass or charge), $a$ undergoes constant acceleration:
$a(t) = \half s_c\,\bar\rho\,t^2$ plus initial conditions.

\subsection{Dirichlet outer wall}

Zero-Dirichlet at the outermost ring forces $a(t) \equiv 0$ at the
wall, which under enforcement at every step pins the DC mode for the
entire box.  The friction term then sees no DC content, and the
static fixed point is the discrete Dirichlet-Poisson solution.  The
\emph{near-wall} value of $\Phi$ is forced to zero rather than
continuing the L-W $\ln r$ tail; the L-W initialization must be
tapered through the absorbing zone to match.

This is the configuration recommended for net-nonzero sources, and
is what the v38 \S 9.6 absorber implicitly relied on (the
multiplicative form does the pinning automatically).

\subsection{Neumann outer wall}

Zero-Neumann ($\partial_n \Phi = 0$ at the wall) leaves the DC mode
unconstrained.  With $\bar\rho \ne 0$ the DC mode runs away
quadratically in time.  For the demo's pic\_binary
($m_\text{total} = 0.02$, $W H = 147456$) at $\bar\rho \sim
10^{-7}$ and $s_c \sim 10$, the DC component grows past $|\Phi^\text{static}|
\sim 10^{-1}$ in $\sim 10^3$ steps -- visually rapid.

Two gauge-fix options:
\begin{enumerate}[leftmargin=*,nosep]
\item \textbf{Zero-mean each step.}  After the leapfrog rotation,
  subtract the spatial mean of $\Phig$ (and $\phi_{em}$) from prev and
  curr.  This is gauge-equivalent (Poisson is defined up to a
  constant; forces depend on gradients only).  Tested implementation:
  works for pic\_static but exhibits a small bulk-momentum drift in
  multi-particle scenarios -- the per-step constant subtraction
  apparently interacts with the centered-friction term in a way
  that bleeds momentum out of the conserved sector.  Cause not yet
  isolated; flagged for follow-up.
\item \textbf{Pin one cell.}  Force $\Phi$ at a single fixed cell to
  zero each step.  Equivalent to a one-point Dirichlet condition;
  preserves Neumann elsewhere.  Not implemented in v42.
\end{enumerate}
The pragmatic choice in v42 is to use \textbf{Dirichlet outer BC}
plus friction.  The friction does the work the multiplicative
damping was supposed to do (absorbing the lowest mode), and the
Dirichlet condition does the work the multiplicative damping was
\emph{also} doing as a side effect (pinning the DC).  Separation of
concerns.

%--------------------------------------------------------------------
\section{Summary of replacement algorithms}
%--------------------------------------------------------------------

\subsection*{Replacement for v38 \S 15.9 (field initialization)}

\begin{algsteps}
  \item Place all particles at their scenario IC positions and momenta.
  \item For each particle, evaluate Eqs.~(\ref{eq:phig_lw})-(\ref{eq:A_lw})
    at every grid cell, \emph{subtracting that particle's boundary-ring
    mean of $\ln r$} from its contribution (gauge shift; all six fields).
    Clip $r$ at $\half\dd x$ at the source.  This leaves the summed field
    $\approx 0$ on the boundary ring, consistent with Dirichlet.  (No
    multiplicative taper -- it distorts gradients.)
  \item Copy the resulting curr buffer into prev and next so the
    encoded time derivative is identically zero.
  \item Settle to the discrete fixed point: freeze the particles and
    iterate the field $\sim\!\max(W,H)$ steps with the friction absorber
    on (wall ramp suffices).  This sheds the near-source
    discrete-vs-continuum residual; the momentum impulse on release
    falls to round-off.  Reset \texttt{step\_count} to 0.
  \item Release the particles and begin dynamics.  No half-step momentum
    back-extrapolation is needed -- the field is already static.
\end{algsteps}

\subsection*{Replacement for v38 \S 9.6 (absorbing boundary layer)}

\begin{algsteps}
  \item Allocate \texttt{friction\_d[k]} of size $W \times H$.
  \item Populate with a \emph{wall-ramp-only} profile: $\beta_\text{floor}
    = 0$ in the interior, rising as $\beta_\text{wall}\,u^2$ over the last
    $N_\text{damp}$ cells from each wall.  Do \emph{not} use a nonzero
    interior floor for dynamics -- it bleeds the orbital near-field energy
    of moving sources (numerical inspiral).
  \item In \texttt{gr\_sim\_step}, after the standard leapfrog
    computes \texttt{next} for every field but \emph{before} the
    three-pointer rotation, apply the per-cell correction
    (\ref{eq:centered_post}).
  \item Outer-wall boundary: zero-Dirichlet (\texttt{next[outer ring]
    = 0}) is recommended for net-nonzero sources.  Neumann is
    available as a flag but requires an additional gauge fix; see
    \S\ref{sec:dc_mode}.
  \item No multiplicative damping is applied -- the entire absorber
    function is carried by the friction.
\end{algsteps}

%--------------------------------------------------------------------
\section{Status \& open items}
\label{sec:status}
%--------------------------------------------------------------------

\paragraph{What works (Revision 2, 2026-06-05).}  The combination
\emph{boundary-mean L-W + frozen-friction settle + Dirichlet +
wall-only friction + BUMP $R=8$ deposition + $v_\text{factor}=1.05$}
gives a stable, near-circular, non-inspiraling, COM-pinned binary:
\begin{itemize}[leftmargin=*,nosep]
  \item No lowest-mode trap, no inversion oscillation, no startup
    transient (boundary-mean + settle).
  \item Center of mass pinned: $\max|\vv{p}_0+\vv{p}_1| \sim 4\times
    10^{-4}\,m v_\text{orb}$ over many orbits.
  \item Separation breathing $\sim 2\%$, secular growth $\approx 1.0$
    (no numerical inspiral), stable through $\sim 8$ orbits.
  \item Discrete static is reached and held.
\end{itemize}

\paragraph{Resolved since Revision 1.}
\begin{itemize}[leftmargin=*,nosep]
  \item \emph{COM drift on pic\_binary} -- the first draft attributed
    this to the L-W startup transient.  Measurement (pre-settling the
    field to its exact fixed point did \emph{not} reduce it) showed the
    dominant cause was instead the \textbf{TSC deposition kernel}.
    \textbf{BUMP $R=8$} pins the COM ($\sim 3700\times$ smaller drift
    than TSC).  Self-field subtraction is \emph{not} needed for this.
  \item \emph{Orbit breathing} -- two compounding causes, both fixed:
    (a) the interior friction floor was bleeding orbital energy (now
    wall-ramp-only friction, \S\ref{sec:friction}); (b) a genuine
    discrete-vs-continuum velocity calibration -- the discrete force
    during motion is $\sim 8\%$ stronger than the continuum
    $\sqrt{G_\text{eff} m}$, so the matched circular speed is
    $v_\text{factor} \approx 1.05$.  Note the \emph{rest}-measured force
    ratio (1.20) overshoots; the moving optimum must be found by
    minimizing breathing, because the velocity-dependent
    (gravitomagnetic / inductive) forces reduce the effective inward
    force during motion.
\end{itemize}

\paragraph{Open follow-ups.}
\begin{itemize}[leftmargin=*,nosep]
  \item \emph{Continuum-discrete Poisson residual.}  See
    \S\ref{sec:lw_discrete}.  Currently shed by the frozen-friction
    settle; could be eliminated at $t=0$ by a TSC-convolved L-W.
  \item \emph{Slow eccentricity creep} of the calibrated binary
    (breathing rises from $\sim 2\%$ at 4--8 orbits to $\sim 5\%$ by
    12).  Minor; likely accumulated discrete-force anisotropy.
  \item \emph{Neumann + zero-mean momentum bleed.}  See
    \S\ref{sec:dc_mode}, item 1.  Worked around by using Dirichlet.
  \item \emph{$v_\text{factor}=1.05$ is empirical} and specific to the
    mass / radius / kernel / grid; a general scenario would need to
    re-calibrate.
\end{itemize}

\end{document}
